\section{Defense Mechanisms}
\label{sec:defenses}

We evaluate five defenses under a \emph{pre-ingestion filtering} model:
before any new entry is stored in memory, the defense is applied; flagged
entries are blocked.  Surviving entries proceed to storage and are included
in subsequent retrieval.  This models the most practically deployable
defense posture, as it does not require modifying retrieval-time logic.

\subsection{Watermark Defense}
\label{sec:watermark_defense}

Content provenance is established at write time using the unigram watermark
of \citet{zhao2024watermark}.  A PRF-keyed partition divides the vocabulary
into \emph{green} (fraction $\gamma=0.25$) and \emph{red} lists; a logit
bias $\delta=2.0$ is applied to green-list items at generation time.
At read time, the z-score
\begin{equation}
  z = \frac{g - \gamma n}{\sqrt{\gamma(1-\gamma)n}}
  \label{eq:z_score}
\end{equation}
is computed ($g$ = green-list count, $n$ = sequence length) and compared to
a threshold $z_{\mathrm{thr}}=4.0$.  Entries with $z < z_{\mathrm{thr}}$ are
flagged as potentially unsigned.

\paragraph{Character-level adaptation.}
The original scheme operates at the LLM token level (GPT-2 vocab, 50{,}257
tokens).  Since memory entries may arrive from sources other than a
watermarked LLM (user input, tool outputs, API responses), we adapt the
scheme to character level: the ASCII printable set (95 characters) is
partitioned via the same PRF, and content generation is biased toward
green-list characters by probabilistic substitution (target green fraction
$\gamma + \delta \times 0.1 = 0.45$).  Character-level detection operates
on Eq.~\ref{eq:z_score} with $n$ = content length in characters.

\paragraph{Token-level scheme (\textsc{Tok-WM}).}
We additionally implement the paper-faithful token-level scheme using GPT-2
as a backbone LLM.  A pre-existing memory entry is re-generated by GPT-2
with green-list logits boosted by $+\delta=2.0$ at every step; the
regenerated continuation is used as the watermarked entry.  The green token
fraction under this bias is:
\begin{equation}
  p_{\text{green}} = \frac{\gamma e^{\delta}}{\gamma e^{\delta} + (1-\gamma)}
  \approx 0.711 \quad (\gamma=0.25,\;\delta=2.0),
\end{equation}
substantially higher than the character-level target of $0.45$.
Detection uses the same z-score (Eq.~\ref{eq:z_score}) computed over GPT-2
tokenization.  The token-level scheme achieves higher statistical power at
short text lengths: for $n=50$ tokens, the expected watermarked z-score is
$\approx 7.5$ vs.\ $\approx 4.1$ for character-level at equivalent length.
Figure~\ref{fig:watermark_comparison} compares the two schemes empirically.

\IfFileExists{figures/fig_p17_watermark_comparison.pdf}{%
\begin{figure}[ht]
  \centering
  \includegraphics[width=0.95\linewidth]{fig_p17_watermark_comparison.pdf}
  \caption{Watermark scheme comparison.
           Left: z-score distributions for clean and watermarked texts
           under character-level and token-level schemes ($n=80$ samples each).
           Right: mean z-score (watermarked) vs.\ text length.
           The token-level scheme provides higher detection power at short
           lengths ($<200$ chars), where the character-level z-score can
           fall below the threshold.
           Both schemes approach $z \gg z_{\text{thr}}$ for longer texts.}
  \label{fig:watermark_comparison}
\end{figure}%
}{}

\subsection{Input Validation Defense}

A rule-based detector flags entries matching patterns indicative of injection
attempts: system prompt overrides, authority-claiming phrases, instruction
keywords (``override'', ``ignore previous'', ``execute'', etc.).  Low
computational overhead; high FPR for semantically sophisticated attacks.

\subsection{Proactive Defense}

Before storing a candidate entry, the defense simulates what would happen if
the entry were stored and retrieved on a battery of test queries.  Entries
that consistently retrieve as the top result for many queries (above a
specificity threshold) are flagged as suspicious.  This is analogous to
\citet{gao2019strip}'s perturbation-based detection.

\subsection{Composite Defense}

A weighted ensemble of Watermark, Validation, and Proactive detectors.
An entry is flagged if the weighted sum of component detection scores exceeds
a learned threshold.  Achieves the highest $\tpr$ and lowest $\fpr$ of the
single-round defenses.

\subsection{Semantic Anomaly Detection (\sad{})}
\label{sec:sad}

\paragraph{Motivation.}
\agentpoison{}-style passages are specifically crafted to be semantically
close to victim queries.  This closeness, necessary for the attack to succeed,
is a detectable signal: \emph{legitimate memory entries from benign sources
are not expected to be highly similar to observed victim queries}.

\paragraph{Algorithm.}
\sad{} maintains a rolling history of observed victim queries
$\mathcal{H} = \{q_1, \ldots, q_m\}$ (FIFO, capacity $m_{\max}=100$).
After calibrating on a benign reference corpus, it computes for each
candidate entry $c$:
\begin{equation}
  s(c) = \max_{q \in \mathcal{H}} \cos(E(c), E(q)).
  \label{eq:sad_score}
\end{equation}
The calibration step estimates $\mu$ and $\sigma$ of $s(\cdot)$ over benign
entries and observed queries.  An entry is flagged as anomalous if:
\begin{equation}
  s(c) > \mu + k \cdot \sigma, \qquad k = 2.0\;\text{(default)}.
  \label{eq:sad_threshold}
\end{equation}

\paragraph{Properties.}
\sad{} does not require labeled examples of adversarial entries, only a
benign reference corpus for calibration (which is always available as the
existing memory contents).  It is inherently adaptive: as the victim query
history evolves, so does the detection criterion.  Batch encoding via
\texttt{detect\_batch()} amortizes encoder inference over all candidate
entries in a single forward pass.

\paragraph{Limitations.}
\sad{} is less effective against \injecmem{}, whose broad-anchor passages are
specifically designed \emph{not} to be suspiciously similar to any particular
query.  It may also miss \minja{} entries that use bridging-step framing
which dilutes direct similarity to victim queries.

\subsection{Adaptive Adversary Analysis}
\label{sec:adaptive}

A key security question for any anomaly-based defense is: \emph{can an
adaptive adversary who knows the detection statistics craft passages that
evade detection while preserving retrieval effectiveness?}

\paragraph{Adaptive attack formulation.}
We model a white-box adversary who has access to the calibration statistics
$\mu$ and $\sigma$ (either through leakage or independent estimation) and
applies greedy synonym substitution to craft passages just below the detection
threshold $\mu + k\sigma$.  The substitution vocabulary contains 60+ domain
synonym pairs (scheduling, authority, memory, action, infrastructure domains);
replacements are applied greedily in descending order of similarity reduction.
Formally, the adversary solves:
\begin{equation}
  \text{minimise}_{w \in \text{Vocab}(c)}\;
  s(\text{substitute}(c, w))
  \quad\text{subject to}\quad
  \text{rank}(c, \mathcal{M}) \leq k,
  \label{eq:adaptive_obj}
\end{equation}
where the rank constraint enforces that the passage still appears in the
top-$k$ retrieval results for at least one victim query.

\paragraph{Fundamental tension.}
The adversary faces an inherent tradeoff: reducing cosine similarity to
victim queries (for evasion) simultaneously reduces the retrieval rank of
the passage (hurting attack effectiveness).  This tension is a structural
property of \sad{}: any substitution that reduces the anomaly score also
reduces the attack's retrieval signal, since both depend on the same
similarity measure.  Formally, the two objectives are in direct conflict:
\begin{equation}
  \nabla_c \bigl[-s(c)\bigr]
  \;=\;
  -\nabla_c \bigl[\max_{q \in \mathcal{H}} \cos(E(c), E(q))\bigr],
\end{equation}
which is also the negative gradient of the retrieval objective for the
highest-similarity victim query.  Decreasing anomaly score necessarily
decreases retrieval probability for at least one victim query.

\paragraph{Empirical results.}
Table~\ref{tab:adaptive_sad} (Section~\ref{sec:adaptive_results}) reports
the evasion rate and retrieval degradation under the adaptive adversary
for all three attacks.  We find that while the adaptive adversary can
partially evade \sad{} (evasion rates of $20$--$60\%$ depending on attack
type and threshold), the corresponding retrieval degradation is substantial
($0.15$--$0.45$ in ASR-R), confirming the fundamental tension empirically.
